\clearpage
\section{The difference $G-A$}\label{sec:GA}

This section and the next one use only the function $h$ and the functions defined from it in the previous sections. The Riemann zeta function is not used. For real $u>0$ with $u\notin\mathbb Z$ put
\begin{equation}\label{eq:ell}
\ell(u)=\log\frac{\pi u}{|\sin(\pi u)|} .
\end{equation}
Since $|\sin(\pi u)|<\pi u$ for $u>0$, we have $\ell(u)>0$. By \CGeqref{eq:h}, $h(u)-\log(2\pi u)=-\ell(u)$, so \CGeqref{eq:GminusA} reads
\begin{equation}\label{eq:GAell}
G(s)-A(s)=-\int_0^\infty\ell(u)\,u^{-s-1}\,du\qquad(0<\operatorname{Re}s<2).
\end{equation}

\begin{proposition}[Sign and size of $G-A$]\label{prop:GAbound}\VCGreviewed{1}
Let $0<\sigma<2$. Then $G(\sigma)-A(\sigma)<0$, and for every real $t$
\begin{equation}\label{eq:GAbound}
|G(\sigma+it)-A(\sigma+it)|\le A(\sigma)-G(\sigma),
\end{equation}
with equality only for $t=0$. In particular $G(s)\ne A(s)$ for real $s$ with $0<s<2$.
\end{proposition}

\begin{proof}
The integral in \CGeqref{eq:GAell} converges absolutely, as stated after \CGeqref{eq:A}. The substitution $u=e^v$ gives $G(\sigma+it)-A(\sigma+it)=-\int_{\mathbb R}f(v)e^{-itv}\,dv$ with $f(v)=\ell(e^v)e^{-\sigma v}$, and $f(v)>0$ for every $v$ with $e^v\notin\mathbb Z$. Hence $G(\sigma)-A(\sigma)=-\int_{\mathbb R}f<0$ and $|G(\sigma+it)-A(\sigma+it)|\le\int_{\mathbb R}f=A(\sigma)-G(\sigma)$. Suppose that equality holds, and choose a real $\theta$ with $\int_{\mathbb R}f(v)e^{-itv}dv=e^{i\theta}\int_{\mathbb R}f$. Then $\int_{\mathbb R}f(v)\bigl(1-\cos(tv+\theta)\bigr)dv=0$. The integrand is nonnegative and $f>0$ almost everywhere, so $\cos(tv+\theta)=1$ for almost every real $v$, which forces $t=0$.
\end{proof}

For $0<u<1$ put
\begin{equation}\label{eq:q0}
q_0(u)=\log\frac{\pi u\,(1-u^2)}{\sin(\pi u)}=\sum_{k\ge1}\frac{\Sigma_{2k}-1}{k}\,u^{2k},\qquad Q_0(s)=\int_0^1q_0(u)\,u^{-s-1}\,du\quad(\operatorname{Re}s<2).
\end{equation}
The series follows from \CGeqref{eq:logsinc} and $\log(1-u^2)=-\sum_{k\ge1}u^{2k}/k$. Its terms are nonnegative, and $\Sigma_{2k}-1=\sum_{m\ge2}m^{-2k}\le4^{-k}+\int_2^\infty\lambda^{-2k}d\lambda=4^{-k}(1+\frac2{2k-1})\le3\cdot4^{-k}$. Hence $0\le q_0(u)\le3\sum_{k\ge1}4^{-k}/k=3\log\frac43$ on $(0,1)$, and termwise integration gives
\begin{equation}\label{eq:Q0series}
Q_0(s)=\sum_{k\ge1}\frac{\Sigma_{2k}-1}{k\,(2k-s)}\qquad(\operatorname{Re}s<2).
\end{equation}
The right side of \CGeqref{eq:Q0series} is meromorphic on $\mathbb C$, with simple poles at the even integers $2k\ge2$, and we use it as the definition of $Q_0$ on $\mathbb C$.

\begin{proposition}[The difference $G-A$ through the digamma function]\label{prop:GAdigamma}\VCGreviewed{1}
For every complex $s\ne0$ that is not an even integer,
\begin{equation}\label{eq:sA}
sA(s)=\log(2\pi)+\frac1s-\psi\Bigl(1-\frac s2\Bigr)-\gamma+s\,Q_0(s)\qquad(\operatorname{Re}s<2),
\end{equation}
and for $0<\operatorname{Re}s<2$, with $Q$ as in \CGeqref{eq:q},
\begin{equation}\label{eq:GAthree}
G(s)-A(s)=\frac{\psi(1-\frac s2)-\psi(1+s)}{s}-Q(s)-Q_0(s),
\end{equation}
where
\begin{equation}\label{eq:digseries}
\frac{\psi(1-\frac s2)-\psi(1+s)}{s}=-\sum_{n\ge1}\frac{3}{(n+s)(2n-s)} .
\end{equation}
For $0<\operatorname{Re}s<2$ the real part of the series in \CGeqref{eq:digseries} is positive, so the real part of the first term of \CGeqref{eq:GAthree} is negative.
\end{proposition}

\begin{proof}
Let $\operatorname{Re}s<2$, $s\ne0$ and $s\notin2\mathbb Z$. For $k\ge1$ we have $1/(k(2k-s))=\frac1s\bigl(1/(k-\frac s2)-1/k\bigr)$. By \CGeqref{eq:psi} with argument $1-\frac s2$, whose real part is positive,
\begin{equation}\label{eq:psihalf}
\sum_{k\ge1}\frac1{k\,(2k-s)}=-\frac1s\sum_{l\ge0}\Bigl(\frac1{l+1}-\frac1{l+1-\frac s2}\Bigr)=-\frac{\psi(1-\frac s2)+\gamma}{s} .
\end{equation}
Writing $\Sigma_{2k}=1+(\Sigma_{2k}-1)$ in \CGeqref{eq:Aseries} and using \CGeqref{eq:psihalf} and \CGeqref{eq:Q0series} gives \CGeqref{eq:sA}. For $0<\operatorname{Re}s<2$, \CGeqref{eq:digamma} gives $G(s)=(\log(2\pi)-\gamma-\psi(s))/s-Q(s)$. Subtracting \CGeqref{eq:sA} divided by $s$, and using $\psi(1+s)=\psi(s)+1/s$, which follows from \CGeqref{eq:psi}, gives \CGeqref{eq:GAthree}. By \CGeqref{eq:psi} again,
\begin{equation}\label{eq:psidiff}
\psi\Bigl(1-\frac s2\Bigr)-\psi(1+s)=\sum_{n\ge1}\Bigl(\frac1{n+s}-\frac1{n-\frac s2}\Bigr)=-\sum_{n\ge1}\frac{3s}{(n+s)(2n-s)},
\end{equation}
which is \CGeqref{eq:digseries}. Write $s=\sigma+it$. Then $\operatorname{Re}\bigl((n+s)(2n-s)\bigr)=2n^2+n\sigma-\sigma^2+t^2\ge2+\sigma-\sigma^2+t^2=(2-\sigma)(1+\sigma)+t^2>0$ for $n\ge1$ and $0<\sigma<2$. Since $\operatorname{Re}(1/w)=\operatorname{Re}w/|w|^2$ for $w\ne0$, every term of the series has positive real part.
\end{proof}

The next lemma and theorem give an analytic proof that $A$ does not vanish in the closed strip $0\le\operatorname{Re}s\le1$ above height $9$.

\begin{lemma}[A lower bound for the real part of the digamma function]\label{lem:repsi}\VCGreviewed{1}
For real $x$ with $\frac12\le x\le1$ and real $y$,
\begin{equation}\label{eq:repsi}
\operatorname{Re}\psi(x+iy)\ge\frac12\log(x^2+y^2)-\gamma-\log2+\frac{y^2}{2x\,(x^2+y^2)} .
\end{equation}
\end{lemma}

\begin{proof}
By \CGeqref{eq:psi}, $\operatorname{Re}\psi(x+iy)-\psi(x)=\sum_{l\ge0}f(l)$ with $f(\lambda)=y^2/((\lambda+x)((\lambda+x)^2+y^2))$ for $\lambda\ge0$. With $w=\lambda+x>0$ one finds
\begin{equation}\label{eq:convex}
\frac{d^2}{dw^2}\,\frac1{w\,(w^2+y^2)}=\frac{12w^4+6w^2y^2+2y^4}{w^3(w^2+y^2)^3}>0,
\end{equation}
so $f$ is convex. For a convex function the chord lies above the graph, so $\int_l^{l+1}f\le\frac12(f(l)+f(l+1))$ for every integer $l\ge0$. Summing over $l$, and using $f(\lambda)\to0$ as $\lambda\to\infty$, gives $\int_0^\infty f\le\sum_{l\ge0}f(l)-\frac12f(0)$. Since $\int_0^\infty f=\int_x^\infty y^2\,dw/(w(w^2+y^2))=\frac12\log(1+y^2/x^2)$, we obtain $\operatorname{Re}\psi(x+iy)\ge\psi(x)+\frac12\log(1+y^2/x^2)+y^2/(2x(x^2+y^2))$. Termwise differentiation of \CGeqref{eq:psi} for real $x>0$ gives $\frac{d}{dx}(\psi(x)-\log x)=\sum_{l\ge0}(l+x)^{-2}-\frac1x>0$, because $(l+x)^{-2}>\int_l^{l+1}(\lambda+x)^{-2}d\lambda$ and the integrals add up to $\frac1x$. By \CGeqref{eq:psi}, $\psi(\frac12)=-\gamma-2\sum_{l\ge0}\bigl(\frac1{2l+1}-\frac1{2l+2}\bigr)=-\gamma-2\log2$. Hence for $x\ge\frac12$ we have $\psi(x)\ge\log x+\psi(\frac12)+\log2=\log x-\gamma-\log2$, and $\log x+\frac12\log(1+y^2/x^2)=\frac12\log(x^2+y^2)$ gives \CGeqref{eq:repsi}.
\end{proof}

\begin{theorem}[$A$ does not vanish above height $9$]\label{thm:Anonzero}\VCGreviewed{1}
Let $0\le\sigma\le1$ and $|t|\ge9$, and put
\begin{equation}\label{eq:Qbound}
\mathcal Q(t)=\log(8\pi)-\frac12\log(1+t^2)-\frac{t^2}{2(t^2+4)}+\frac1{t^2}-0.64\,\frac{t^2-1}{t^2+4}+10^{-27} .
\end{equation}
Then $\operatorname{Re}\bigl(sA(s)\bigr)\le\mathcal Q(t)<0$ for $s=\sigma+it$. In particular $A(s)\ne0$.
\end{theorem}

\begin{proof}
Since $A(\bar s)=\overline{A(s)}$ by \CGeqref{eq:Aseries}, we may assume $t\ge9$. We bound the real parts of the terms of \CGeqref{eq:sA}. First, $\operatorname{Re}(1/s)=\sigma/(\sigma^2+t^2)\le1/t^2$. Second, apply Lemma~\CGref{lem:repsi} with $x=1-\frac\sigma2\in[\frac12,1]$ and $y=-\frac t2$. Then $x^2+y^2\ge\frac14(1+t^2)$, and $x\le1$, $x^2+y^2\le\frac14(4+t^2)$ give $y^2/(2x(x^2+y^2))\ge t^2/(2(t^2+4))$. Hence
\begin{equation}\label{eq:psibound}
-\operatorname{Re}\psi\Bigl(1-\frac s2\Bigr)-\gamma\le2\log2-\frac12\log(1+t^2)-\frac{t^2}{2(t^2+4)} .
\end{equation}
Third, by \CGeqref{eq:Q0series}, $\operatorname{Re}(sQ_0(s))=\sum_{k\ge1}\frac{\Sigma_{2k}-1}k\,\mathfrak r_k$ with
\begin{equation}\label{eq:rhok}
\mathfrak r_k=\operatorname{Re}\frac{s}{2k-s}=\frac{2k\sigma-\sigma^2-t^2}{(2k-\sigma)^2+t^2} .
\end{equation}
For $0\le\sigma\le1$ we have $2k\sigma-\sigma^2\le2k-1$, because the left side increases with $\sigma$ on $[0,1]$. For $k=1$ the numerator is at most $1-t^2<0$ and the denominator is at most $4+t^2$, so $\rho_1\le-(t^2-1)/(t^2+4)$. If $2\le k$ and $2k-1<t^2$, then $\mathfrak r_k<0$. If $2k-1\ge t^2$, then $k\ge41$ because $t^2\ge81$, and $\mathfrak r_k\le(2k-1)/(2k-1)^2=1/(2k-1)$, since the denominator is at least $(2k-1)^2$. With $\Sigma_{2k}-1\le3\cdot4^{-k}$ these terms add up to at most $\frac{3}{41\cdot81}\sum_{k\ge41}4^{-k}=4^{-40}/3321<10^{-27}$. Finally $\Sigma_2-1>\sum_{m=2}^{10}m^{-2}+\int_{11}^\infty\lambda^{-2}d\lambda>0.5497+0.0909>0.64$. Adding up, $\operatorname{Re}(sA(s))\le\mathcal Q(t)$.

Each of the terms $-\frac12\log(1+t^2)$, $-\frac{t^2}{2(t^2+4)}=-\frac12+\frac2{t^2+4}$, $\frac1{t^2}$ and $-0.64\frac{t^2-1}{t^2+4}=-0.64+\frac{3.2}{t^2+4}$ decreases for $t>0$, so $\mathcal Q(t)\le\mathcal Q(9)$ for $t\ge9$. Since $e^{3.2242}>25.1328>8\pi$ and $e^{4.4066}<82$, we have $\log(8\pi)<3.2242$ and $\frac12\log82>2.2033$. With $\frac{81}{170}>0.4764$, $\frac1{81}<0.0124$ and $0.64\cdot\frac{80}{85}>0.6023$ this gives $\mathcal Q(9)<3.2242-2.2033-0.4764+0.0124-0.6023+10^{-27}<-0.045$.
\end{proof}

\begin{remark}[A computer-assisted complement]\label{rem:Acert}
Theorem~\CGref{thm:Anonzero} leaves the rectangle $0\le\operatorname{Re}s\le1$, $|\operatorname{Im}s|<9$, and its proof gives nothing there. On the real segment $0<s<2$ we have $A(s)>0$ by Proposition~\CGref{prop:A}. The program \texttt{A\_nonvanishing\_check.py} of the files that accompany this note encloses $s^2A(s)$, written through \CGeqref{eq:sA}, in ball arithmetic \cite{Johansson2017} on $188892$ boxes that cover $0\le\operatorname{Re}s\le1$, $0\le\operatorname{Im}s\le102$, and none of the enclosures contains $0$. This is a computer-assisted verification of $A(s)\ne0$ on that rectangle. No other statement of this note depends on it.
\end{remark}

\section{A positive convex kernel for $G-A$}\label{sec:B}

Put $\mathfrak b=\log(2\pi)-\frac32$. Since $2\pi>e^{3/2}$, we have $\mathfrak b>0$. For real $u\ge0$ put
\begin{equation}\label{eq:Bdef}
\mathcal B(u)=4\pi^2\int_0^u(u-\lambda)\,\ell(\lambda)\,d\lambda,\qquad\Omega(s)=\int_0^\infty\mathcal B(u)\,u^{-s-3}\,du\quad(0<\operatorname{Re}s<2).
\end{equation}
The integral defining $\mathcal B$ converges because $\ell$ is nonnegative and locally integrable on $[0,\infty)$: it has only logarithmic singularities at the positive integers, and $\ell(\lambda)\le\pi^2\lambda^2/3$ for $0<\lambda\le\frac12$ by the bound $|\log(\sin(\pi\tau)/(\pi\tau))|\le\pi^2\tau^2/3$ from the proof of Lemma~\CGref{lem:Klocal}.

\begin{proposition}[The kernel $\mathcal B$]\label{prop:B}\VCGreviewed{1}
The function $\mathcal B$ has the following properties.
\textup{(i)} $\mathcal B$ is continuously differentiable on $[0,\infty)$, $\mathcal B'(u)=4\pi^2\int_0^u\ell$ is absolutely continuous and nondecreasing, $\mathcal B''=4\pi^2\ell$ almost everywhere, $\mathcal B$ is convex, and $\mathcal B(u)>0$ for $u>0$.
\textup{(ii)} For every $u>0$,
\begin{equation}\label{eq:BK}
\mathcal B(u)=\mathcal K(u)+2\pi^2u^2(\log u+\mathfrak b).
\end{equation}
In particular $\mathcal B(1)=2\pi^2\mathfrak b$ and $\mathcal B'(1)=2\pi^2(1+2\mathfrak b)$.
\textup{(iii)} For $0\le u\le1$,
\begin{equation}\label{eq:Bseries}
\mathcal B(u)=4\pi^2\sum_{k\ge1}\frac{\Sigma_{2k}}{k\,(2k+1)(2k+2)}\,u^{2k+2},
\end{equation}
and $0<\mathcal B(u)<2\pi^2\mathfrak b\,u^4$ for $0<u<1$.
\textup{(iv)} The integral $\Omega(s)$ converges absolutely for $0<\operatorname{Re}s<2$, and there
\begin{equation}\label{eq:GAOmega}
G(s)-A(s)=-\frac{(s+1)(s+2)}{4\pi^2}\,\Omega(s).
\end{equation}
\end{proposition}

\begin{proof}
\textup{(i)} By Fubini's theorem $\int_0^u(u-\lambda)\ell(\lambda)d\lambda=\int_0^u\bigl(\int_0^\mu\ell\bigr)d\mu$. The inner integral is a continuous nondecreasing function of $\mu$, absolutely continuous with derivative $\ell$ almost everywhere. This gives the statements on $\mathcal B'$ and $\mathcal B''$. A function with nondecreasing derivative is convex, and $\mathcal B(u)>0$ for $u>0$ because $\ell>0$.

\textup{(ii)} By \CGeqref{eq:Kprime}, $\mathcal K(0)=0$ and Fubini's theorem, $\mathcal K(u)=\int_0^u\mathcal K'(\mu)\,d\mu=-4\pi^2\int_0^u\int_0^\mu h(\lambda)\,d\lambda\,d\mu=-4\pi^2\int_0^u(u-\lambda)h(\lambda)\,d\lambda$. Since $\ell(\lambda)+h(\lambda)=\log(2\pi\lambda)$ and $\int_0^u(u-\lambda)\log(2\pi\lambda)\,d\lambda=\frac{u^2}2\log(2\pi u)-\frac{3u^2}4$, we get $\mathcal B(u)-\mathcal K(u)=2\pi^2u^2(\log(2\pi u)-\frac32)$, which is \CGeqref{eq:BK}. At $u=1$, \CGeqref{eq:Kprops} gives $\mathcal K(1)=\mathcal K'(1)=0$, hence $\mathcal B(1)=2\pi^2\mathfrak b$ and $\mathcal B'(1)=2\pi^2(2\cdot1\cdot\log1+1+2\mathfrak b)=2\pi^2(1+2\mathfrak b)$.

\textup{(iii)} For $0\le\lambda<1$, \CGeqref{eq:logsinc} gives $\ell(\lambda)=\sum_{k\ge1}\Sigma_{2k}\lambda^{2k}/k$ with nonnegative terms. Since $\int_0^u(u-\lambda)\lambda^{2k}d\lambda=u^{2k+2}/((2k+1)(2k+2))$, monotone convergence gives \CGeqref{eq:Bseries} for $0\le u\le1$. All the coefficients are positive and all the exponents are at least $4$. Hence for $0<u<1$ every term with $k\ge2$ is strictly smaller than its value at $u=1$ times $u^4$, the term with $k=1$ equals its value at $u=1$ times $u^4$, and $0<\mathcal B(u)<u^4\mathcal B(1)=2\pi^2\mathfrak b\,u^4$.

\textup{(iv)} Let $\sigma=\operatorname{Re}s$. Near $0$ we have $\mathcal B(u)u^{-\sigma-3}<2\pi^2\mathfrak b\,u^{1-\sigma}$, which is integrable for $\sigma<2$. For $u\ge1$, \CGeqref{eq:Kprops} and \CGeqref{eq:BK} give $\mathcal B(u)\le\frac52+2\pi^2u^2(\log u+\mathfrak b)$, and $\mathcal B(u)u^{-\sigma-3}$ is integrable on $[1,\infty)$ for $\sigma>0$. Let $0<\tau<\frac12<X$. Since $\mathcal B'$ is absolutely continuous with $\mathcal B''=4\pi^2\ell$, two integrations by parts give
\begin{equation}\label{eq:Bparts}
4\pi^2\int_\tau^X\ell(u)\,u^{-s-1}\,du=\Bigl[\mathcal B'(u)u^{-s-1}+(s+1)\mathcal B(u)u^{-s-2}\Bigr]_\tau^X+(s+1)(s+2)\int_\tau^X\mathcal B(u)\,u^{-s-3}\,du .
\end{equation}
For $u\le\frac12$ the bound $\ell(\lambda)\le\pi^2\lambda^2/3$ gives $\mathcal B'(u)\le4\pi^4u^3/9$ and $\mathcal B(u)\le\pi^4u^4/9$, so the terms at $\tau$ are $O(\tau^{2-\sigma})$ and tend to $0$. The series of derivatives in the proof of Proposition~\CGref{prop:K} gives $|\mathcal K'|\le2\pi\Sigma_2$, so by \CGeqref{eq:BK} $\mathcal B'(X)=O(X\log X)$ and $\mathcal B(X)=O(X^2\log X)$, and the terms at $X$ are $O(X^{-\sigma}\log X)$ and tend to $0$. Letting $\tau\to0$ and $X\to\infty$ and using \CGeqref{eq:GAell} gives \CGeqref{eq:GAOmega}.
\end{proof}

By \CGeqref{eq:GAOmega}, and since $(s+1)(s+2)\ne0$ for $0<\operatorname{Re}s<2$, the solutions of $G(s)=A(s)$ in the strip $0<\operatorname{Re}s<2$ are exactly the zeros of $\Omega$ there.

\begin{proposition}[An inversion inequality and two positive weights]\label{prop:Binv}\VCGreviewed{1}
For $u>1$ and $0<\sigma\le1$,
\begin{equation}\label{eq:Binv}
\mathcal B(u)>u^{2\sigma+4}\,\mathcal B(1/u) .
\end{equation}
For real $\sigma$ with $0<\sigma<2$ and real $v$ put
\begin{equation}\label{eq:upsilon}
\upsilon_\sigma(v)=e^{-(\sigma+2)v}\mathcal B(e^v),\qquad\upsilon_\sigma^\pm(v)=\upsilon_\sigma(v)\pm\upsilon_\sigma(-v)\quad(v>0).
\end{equation}
Then $\Omega(\sigma+it)=\int_{\mathbb R}\upsilon_\sigma(v)e^{-itv}\,dv$ for real $t$, and
\begin{equation}\label{eq:ReImOmega}
\operatorname{Re}\Omega(\sigma+it)=\int_0^\infty\upsilon^+_\sigma(v)\cos(tv)\,dv,\qquad\operatorname{Im}\Omega(\sigma+it)=-\int_0^\infty\upsilon^-_\sigma(v)\sin(tv)\,dv .
\end{equation}
For $0<\sigma\le1$ and $v>0$ we have $\upsilon^+_\sigma(v)>\upsilon^-_\sigma(v)>0$. The function $\Omega$ is positive on the real interval $(0,2)$ and strictly decreasing on $(0,1]$.
\end{proposition}

\begin{proof}
Let $u>1$ and $0<\sigma\le1$. By Proposition~\CGref{prop:B}\textup{(iii)} with $1/u$ in place of $u$, $u^{2\sigma+4}\mathcal B(1/u)<2\pi^2\mathfrak b\,u^{2\sigma}\le2\pi^2\mathfrak b\,u^2$. By \CGeqref{eq:BK} and $\mathcal K\ge0$, $\mathcal B(u)\ge2\pi^2u^2(\log u+\mathfrak b)>2\pi^2\mathfrak b\,u^2$. This proves \CGeqref{eq:Binv}. The substitution $u=e^v$ in \CGeqref{eq:Bdef} gives the integral over $\mathbb R$. Splitting it at $v=0$ and replacing $v$ by $-v$ on $(-\infty,0)$ gives \CGeqref{eq:ReImOmega}. For $v>0$, \CGeqref{eq:Binv} with $u=e^v$ is $\upsilon_\sigma(v)>\upsilon_\sigma(-v)$, and $\upsilon_\sigma(-v)>0$ by Proposition~\CGref{prop:B}\textup{(i)}. For real $\sigma\in(0,2)$, $\Omega(\sigma)=\int_{\mathbb R}\upsilon_\sigma>0$. On a compact interval $[\sigma_1,\sigma_2]\subset(0,2)$ we have $|v|\upsilon_\sigma(v)\le|v|(\upsilon_{\sigma_1}(v)+\upsilon_{\sigma_2}(v))$, which is integrable, so differentiation under the integral sign is allowed and $\Omega'(\sigma)=-\int_{\mathbb R}v\,\upsilon_\sigma(v)\,dv=-\int_0^\infty v\,\upsilon^-_\sigma(v)\,dv$. For $0<\sigma\le1$ this is negative.
\end{proof}

\begin{theorem}[A region without solutions of $G=A$ near the real axis]\label{thm:GAexcl}\VCGreviewed{1}
For $0<\sigma<2$ and $j=0,1,2$ put $\mathcal M_j(\sigma)=\int_0^\infty(\log u)^j\,\mathcal B(u)\,u^{-\sigma-3}\,du$. If $t$ is real and
\begin{equation}\label{eq:GAexcl}
t^2\bigl(\mathcal M_0(\sigma)\mathcal M_2(\sigma)-\mathcal M_1(\sigma)^2\bigr)<2\,\mathcal M_0(\sigma)^2,
\end{equation}
then $G(\sigma+it)\ne A(\sigma+it)$. In particular $G(\sigma+it)\ne A(\sigma+it)$ whenever $|t|<\sqrt{2\mathcal M_0(\sigma)/\mathcal M_2(\sigma)}$.
\end{theorem}

\begin{proof}
By the substitution $u=e^v$, $\mathcal M_j(\sigma)=\int_{\mathbb R}v^j\upsilon_\sigma(v)\,dv$, and these integrals converge by Proposition~\CGref{prop:B}\textup{(iii)} and \textup{(iv)}. Put $\mu=\mathcal M_1/\mathcal M_0$. By Proposition~\CGref{prop:Binv}, $e^{it\mu}\Omega(\sigma+it)=\int_{\mathbb R}\upsilon_\sigma(v)e^{-it(v-\mu)}dv$. Its real part is $\int_{\mathbb R}\upsilon_\sigma(v)\cos(t(v-\mu))\,dv$, and $\cos y\ge1-\frac{y^2}2$ for real $y$ gives
\begin{equation}\label{eq:Omegacos}
|\Omega(\sigma+it)|\ge\mathcal M_0-\frac{t^2}2\int_{\mathbb R}(v-\mu)^2\upsilon_\sigma(v)\,dv=\mathcal M_0-\frac{t^2}2\Bigl(\mathcal M_2-\frac{\mathcal M_1^2}{\mathcal M_0}\Bigr).
\end{equation}
Under \CGeqref{eq:GAexcl} the right side is positive, and \CGeqref{eq:GAOmega} gives $G\ne A$. The last statement follows because $\mathcal M_0\mathcal M_2-\mathcal M_1^2\le\mathcal M_0\mathcal M_2$.
\end{proof}

\begin{proposition}[Decay of $\Omega$]\label{prop:Omegadecay}\VCGreviewed{1}
Let $0<\sigma<2$. The function $\upsilon_\sigma$ is continuously differentiable on $\mathbb R$, $\upsilon_\sigma'$ is absolutely continuous, $\upsilon_\sigma''$ is integrable on $\mathbb R$, and for real $t\ne0$
\begin{equation}\label{eq:Omegadecay}
|\Omega(\sigma+it)|\le\frac1{t^2}\int_{\mathbb R}|\upsilon_\sigma''(v)|\,dv .
\end{equation}
\end{proposition}

\begin{proof}
Put $u=e^v$. By Proposition~\CGref{prop:B}\textup{(i)}, $\upsilon_\sigma'(v)=u^{-\sigma-2}\bigl(u\mathcal B'(u)-(\sigma+2)\mathcal B(u)\bigr)$ is continuous and absolutely continuous, and almost everywhere $\upsilon_\sigma''(v)=u^{-\sigma-2}\bigl(4\pi^2u^2\ell(u)-(2\sigma+3)u\mathcal B'(u)+(\sigma+2)^2\mathcal B(u)\bigr)$. With $dv=du/u$, the three terms contribute to $\int_{\mathbb R}|\upsilon_\sigma''|$ at most constant multiples of $\int_0^\infty\ell(u)u^{-\sigma-1}du$, $\int_0^\infty\mathcal B'(u)u^{-\sigma-2}du$ and $\int_0^\infty\mathcal B(u)u^{-\sigma-3}du$. The first is finite by \CGeqref{eq:GAell}, and the other two by the bounds for $\mathcal B$ and $\mathcal B'$ near $0$ and near $\infty$ in the proof of Proposition~\CGref{prop:B}\textup{(iv)}. The same bounds show that $\upsilon_\sigma(v)$ and $\upsilon_\sigma'(v)$ tend to $0$ as $v\to\pm\infty$. Two integrations by parts in $\Omega(\sigma+it)=\int_{\mathbb R}\upsilon_\sigma(v)e^{-itv}dv$ give $\Omega(\sigma+it)=-t^{-2}\int_{\mathbb R}\upsilon_\sigma''(v)e^{-itv}dv$, and \CGeqref{eq:Omegadecay} follows.
\end{proof}

The inequality \CGeqref{eq:Binv} holds for every $\sigma\in(0,1]$, so Proposition~\CGref{prop:Binv} does not distinguish the value $\sigma=\frac12$. The next proposition shows that positivity and periodicity of the kernel alone cannot place the solutions of $G=A$ on $\operatorname{Re}s=\frac12$. For real $x$ put
\begin{equation}\label{eq:Ksharp}
\mathcal K^\sharp(x)=1-\cos(2\pi x)+2^{-11/4}\bigl(1-\cos(4\pi x)\bigr) .
\end{equation}
Like $\mathcal K$, the function $\mathcal K^\sharp$ is even, has period $1$, satisfies $\mathcal K^\sharp>0$ off the integers and $\mathcal K^\sharp(n)=(\mathcal K^\sharp)'(n)=0$ at the integers $n$. Put $H^\sharp(s)=\int_1^\infty\mathcal K^\sharp(u)u^{-s-3}du$ for $\operatorname{Re}s>-2$, $\widetilde H^\sharp(s)=\int_0^1\mathcal K^\sharp(u)u^{-s-3}du$ for $\operatorname{Re}s<0$, and
\begin{equation}\label{eq:Cdef}
\mathcal C(s)=\int_0^\infty(1-\cos y)\,y^{-s-3}\,dy\qquad(-2<\operatorname{Re}s<0).
\end{equation}
In view of \CGeqref{eq:factor} and \CGeqref{eq:Afactor}, the counterparts of $G$ and $A$ for $\mathcal K^\sharp$ are $G^\sharp=-\frac{(s+1)(s+2)}{4\pi^2}H^\sharp$ and $A^\sharp=\frac{(s+1)(s+2)}{4\pi^2}\widetilde H^\sharp$.

\begin{proposition}[A positive kernel with solutions off the line]\label{prop:Ksharp}\VCGreviewed{1}
The functions $\widetilde H^\sharp$ and $\mathcal C$ extend to meromorphic functions on $\operatorname{Re}s>-2$ whose poles are real, and on $\operatorname{Re}s>-2$
\begin{equation}\label{eq:GAsharp}
G^\sharp(s)-A^\sharp(s)=-\frac{(s+1)(s+2)}{4\pi^2}\,(2\pi)^{s+2}\bigl(1+2^{s-3/4}\bigr)\,\mathcal C(s) .
\end{equation}
Consequently $G^\sharp(s)=A^\sharp(s)$ at the points $s=\frac34+i(2k+1)\pi/\log2$, $k\in\mathbb Z$, which have real part $\frac34$ and nonzero imaginary part.
\end{proposition}

\begin{proof}
The Taylor series $\mathcal K^\sharp(u)=\sum_{n\ge1}\mathfrak c_nu^{2n}$ at $0$ converges for all $u$, with $|\mathfrak c_n|\le(4\pi)^{2n}/(2n)!$. Termwise integration gives $\widetilde H^\sharp(s)=\sum_{n\ge1}\mathfrak c_n/(2n-2-s)$ for $\operatorname{Re}s<0$, and the right side is meromorphic on $\mathbb C$ with poles at the even integers $2n-2\ge0$. In the same way $\int_0^1(1-\cos y)y^{-s-3}dy=\sum_{n\ge1}(-1)^{n+1}/((2n)!(2n-2-s))$, while $\int_1^\infty(1-\cos y)y^{-s-3}dy$ and $H^\sharp(s)$ are holomorphic on $\operatorname{Re}s>-2$ because $1-\cos y$ and $\mathcal K^\sharp$ are bounded. Let $-2<\operatorname{Re}s<0$. The substitutions $y=2\pi u$ and $y=4\pi u$ give
\begin{equation}\label{eq:HHsharp}
H^\sharp(s)+\widetilde H^\sharp(s)=\int_0^\infty\mathcal K^\sharp(u)\,u^{-s-3}\,du=\bigl((2\pi)^{s+2}+2^{-11/4}(4\pi)^{s+2}\bigr)\mathcal C(s)=(2\pi)^{s+2}\bigl(1+2^{s-3/4}\bigr)\mathcal C(s),
\end{equation}
and multiplication by $-\frac{(s+1)(s+2)}{4\pi^2}$ gives \CGeqref{eq:GAsharp} there. Both sides of \CGeqref{eq:GAsharp} are meromorphic on the connected set $\operatorname{Re}s>-2$, so the identity holds on all of it. At $s_k=\frac34+i(2k+1)\pi/\log2$ we have $2^{s_k-3/4}=e^{i(2k+1)\pi}=-1$, and $s_k$ is not real, so it is not a pole of $\mathcal C$ or of $A^\sharp$. Hence the right side of \CGeqref{eq:GAsharp} vanishes at $s_k$.
\end{proof}

\section{The difference $G-A$ and the Riemann zeta function}\label{sec:zeta}

In this section, and only here and in Section~\CGref{sec:fig}, $\zeta$ denotes the Riemann zeta function, and $\Gamma$ denotes Euler's gamma function. We use three cited results. The functional equation in the form \cite[(2.1.8), p.~16]{Titchmarsh1986}
\begin{equation}\label{eq:zetafe}
\zeta(1-s)=2^{1-s}\pi^{-s}\cos\Bigl(\frac{\pi s}2\Bigr)\Gamma(s)\zeta(s) .
\end{equation}
The evaluation $\int_0^\infty y^{-w-1}\sin y\,dy=-\Gamma(-w)\sin(\frac{\pi w}2)$ for $-1<\operatorname{Re}w<0$, which is the last step of the first proof of the functional equation in \cite[p.~15]{Titchmarsh1986}. The relations $\Gamma(z+1)=z\Gamma(z)$ and $\Gamma(z)\Gamma(1-z)=\pi/\sin(\pi z)$ for non-integral $z$, and the fact that $1/\Gamma$ is an entire function \cite[\S4.41, p.~149]{Titchmarsh1939}.

\begin{proposition}[The difference $G-A$ through $\zeta$]\label{prop:zeta}\VCGreviewed{1}
For $-2<\operatorname{Re}s<0$,
\begin{equation}\label{eq:HHzeta}
H(s)+\widetilde H(s)=\int_0^\infty\mathcal K(u)\,u^{-s-3}\,du=(2\pi)^{s+2}\zeta(1-s)\,\mathcal C(s),
\end{equation}
and on $\mathbb C$, as an identity between meromorphic functions,
\begin{equation}\label{eq:GAzeta}
G(s)-A(s)=\frac\pi s\cot\Bigl(\frac{\pi s}2\Bigr)\zeta(s) .
\end{equation}
\end{proposition}

\begin{proof}
Let $-2<\sigma<0$. The terms of \CGeqref{eq:K} are nonnegative, and the substitution $y=2\pi mu$ gives $\int_0^\infty(1-\cos(2\pi mu))u^{-\sigma-3}du=(2\pi m)^{\sigma+2}\mathcal C(\sigma)$. Since $\sum_{m\ge1}m^{-3}(2\pi m)^{\sigma+2}=(2\pi)^{\sigma+2}\sum_{m\ge1}m^{\sigma-1}<\infty$, the theorems of Tonelli and Fubini allow termwise integration for complex $s$ with $\operatorname{Re}s=\sigma$, and $\sum_{m\ge1}m^{s-1}=\zeta(1-s)$. This gives \CGeqref{eq:HHzeta}, where $H$ and $\widetilde H$ are as in \CGeqref{eq:H} and \CGeqref{eq:Afactor}. Now let $-2<\sigma<-1$. Integration by parts in \CGeqref{eq:Cdef}, whose boundary terms vanish because $(1-\cos y)y^{-\sigma-2}$ tends to $0$ at $0$ and at $\infty$, and the evaluation quoted above with $w=s+1$ give
\begin{equation}\label{eq:Cgamma}
\mathcal C(s)=\frac1{s+2}\int_0^\infty y^{-s-2}\sin y\,dy=-\frac{\Gamma(-s-1)}{s+2}\sin\Bigl(\frac{\pi(s+1)}2\Bigr)=-\frac{\Gamma(-s-1)}{s+2}\cos\Bigl(\frac{\pi s}2\Bigr).
\end{equation}
By \CGeqref{eq:factor}, valid on $\mathbb C$ by Theorem~\CGref{thm:entire}, and by \CGeqref{eq:Afactor}, $G-A=-\frac{(s+1)(s+2)}{4\pi^2}(H+\widetilde H)$ for $-2<\sigma<0$. With \CGeqref{eq:HHzeta} and $(s+1)\Gamma(-s-1)=-\Gamma(-s)$ this gives, for $-2<\sigma<-1$,
\begin{equation}\label{eq:GAgamma}
G(s)-A(s)=-(2\pi)^s\,\Gamma(-s)\cos\Bigl(\frac{\pi s}2\Bigr)\zeta(1-s)=-2\,\Gamma(s)\Gamma(-s)\cos^2\Bigl(\frac{\pi s}2\Bigr)\zeta(s),
\end{equation}
by \CGeqref{eq:zetafe}. Since $\Gamma(1-s)=-s\Gamma(-s)$, we have $\Gamma(s)\Gamma(-s)=-\pi/(s\sin(\pi s))$, and $\sin(\pi s)=2\sin(\frac{\pi s}2)\cos(\frac{\pi s}2)$ gives \CGeqref{eq:GAzeta} for $-2<\sigma<-1$. Both sides of \CGeqref{eq:GAzeta} are meromorphic on $\mathbb C$, so the identity holds on $\mathbb C$.
\end{proof}

\begin{corollary}[The condition on $G$ and $A$ is the Riemann hypothesis]\label{cor:RH}\VCGreviewed{1}
For every non-real $s$, $G(s)=A(s)$ if and only if $\zeta(s)=0$, and the multiplicity of $s$ as a zero of $G-A$ equals its multiplicity as a zero of $\zeta$. Hence the following two statements are equivalent.
\textup{(a)} Every $s$ with $\operatorname{Im}s\ne0$ and $G(s)=A(s)$ has the form $s=\frac12+it$ with real $t\ne0$.
\textup{(b)} Every zero of $\zeta$ that is not real has real part $\frac12$, that is, the Riemann hypothesis.
\end{corollary}

\begin{proof}
Let $s$ be non-real. The poles of $A$ are real by Proposition~\CGref{prop:A}, so $G(s)-A(s)$ is finite. The function $\frac\pi s\cot(\frac{\pi s}2)$ is holomorphic and nonzero at $s$, because $\sin z$ vanishes only at the real points $z\in\pi\mathbb Z$ and $\cos z$ only at the real points $z\in\frac\pi2+\pi\mathbb Z$. Hence \CGeqref{eq:GAzeta} gives the first statement, with multiplicities, and the equivalence of \textup{(a)} and \textup{(b)} follows.
\end{proof}

Through Corollary~\CGref{cor:RH}, every theorem on the non-real zeros of $\zeta$ is a theorem on the non-real solutions of $G=A$. For instance, $\zeta(1+it)\ne0$ for real $t\ne0$ \cite[\S3.1, p.~45]{Titchmarsh1986}, so $G(1+it)\ne A(1+it)$ for $t\ne0$, and also $G(1)\ne A(1)$ by Proposition~\CGref{prop:GAbound}. Let $N(T)$ be the number of zeros of $\zeta$ with $0\le\operatorname{Re}s\le1$ and $0<\operatorname{Im}s\le T$, counted with multiplicity \cite[\S9.2, p.~211]{Titchmarsh1986}. By Corollary~\CGref{cor:RH}, $N(T)$ is also the number of solutions of $G=A$ in the same region, and by \cite[Theorem~9.4, p.~214]{Titchmarsh1986}
\begin{equation}\label{eq:NT}
N(T)=\frac T{2\pi}\log\frac T{2\pi e}+O(\log T)\qquad(T\to\infty).
\end{equation}
Put $N_0(T)=\frac T{2\pi}\log\frac T{2\pi e}+\frac78$, the main term of \CGeqref{eq:NT} plus the constant $\frac78$. On $T>2\pi$ the function $N_0$ increases from $-\frac18$ to $\infty$. For every integer $k\ge1$ let $\tau_k>2\pi$ be the solution of $N_0(\tau_k)=k-\frac12$. This equation reads $\frac{\tau_k}{2\pi e}\log\frac{\tau_k}{2\pi e}=(k-\frac{11}8)/e$ with $\log\frac{\tau_k}{2\pi e}>-1$, so that $\log\frac{\tau_k}{2\pi e}=W((k-\frac{11}8)/e)$ and
\begin{equation}\label{eq:tauk}
\tau_k=\frac{2\pi(k-\frac{11}8)}{W\bigl((k-\frac{11}8)/e\bigr)},
\end{equation}
where $W$ is the inverse function of $w\mapsto we^w$ on $[-1,\infty)$. The first values are $\tau_1=14.5213$, $\tau_2=20.6557$, $\tau_3=25.4927$, $\tau_4=29.7394$, $\tau_5=33.6245$ and $\tau_6=37.2574$. The imaginary parts of the six solutions of $G=A$ shown in Figure~\CGref{fig:zeromap}, which are those of the first six non-real zeros of $\zeta$ in the upper half-plane, differ from $\tau_1,\dots,\tau_6$ by at most $0.69$. At the points $\frac12+i\tau_k$ with $k\le6$ the computed values of $|G-A|$ lie between $0.054$ and $0.105$, so these points are not solutions of $G=A$. The statements of the last two sentences are numerical observations.

\section{Conclusions}\label{sec:concl}

Sections~\CGref{sec:GA} and~\CGref{sec:B} characterize the difference $G-A$ without the Riemann zeta function, using only the function $h$. On the strip $0<\operatorname{Re}s<2$ it is minus the Mellin transform \CGeqref{eq:GAell} of the positive function $\ell(u)=\log(\pi u/|\sin(\pi u)|)$, so that $|G-A|$ is bounded by its value $A(\sigma)-G(\sigma)$ on the real axis, which is attained only there. It is also the sum \CGeqref{eq:GAthree} of an explicit digamma term, whose real part is negative, and of two transforms of nonnegative remainders. After the factor $-(s+1)(s+2)/(4\pi^2)$, it is the Mellin transform $\Omega$ of the convex kernel $\mathcal B$ of \CGeqref{eq:Bdef}. This kernel is positive and continuously differentiable, and on $(0,1)$ it is a power series with positive coefficients and exponents at least $4$. It satisfies the strict inversion inequality \CGeqref{eq:Binv}, so that the real and imaginary parts of $\Omega$ are cosine and sine transforms of two positive weights. It follows that $\Omega$ is positive and decreasing on the real segment $(0,1]$, that there is the region \CGeqref{eq:GAexcl} without solutions of $G=A$ around the real axis, and that $\Omega(\sigma+it)=O(t^{-2})$. Theorem~\CGref{thm:Anonzero} proves that $A(s)\ne0$ for $0\le\operatorname{Re}s\le1$ and $|\operatorname{Im}s|\ge9$, so that there the equation $G=A$ is the equation $G/A=1$. The condition on $G$ and $A$ that every solution of $G(s)=A(s)$ with $\operatorname{Im}s\ne0$ has the form $s=\frac12+it$ is a statement about these two functions alone, and it can be formulated without $\zeta$.

None of these properties places the solutions on the line $\operatorname{Re}s=\frac12$. The inversion inequality holds for every $\sigma\in(0,1]$ and does not single out $\sigma=\frac12$. Proposition~\CGref{prop:Ksharp} gives a kernel $\mathcal K^\sharp$ with two harmonics that is nonnegative, has period $1$ and vanishes to second order at the integers, exactly like $\mathcal K$, and for which the counterpart of $G-A$ vanishes at points with real part $\frac34$. Therefore a proof of the condition has to use the coefficients $m^{-3}$ of all the harmonics of $\mathcal K$, and not only the positivity and the periodicity of the kernel.

Proposition~\CGref{prop:zeta} shows where these coefficients go. The sum over all the harmonics of $\mathcal K$ is the Dirichlet series of $\zeta(1-s)$, and $G-A$ equals $\frac\pi s\cot(\frac{\pi s}2)\zeta(s)$. This is the problem that appears when $\zeta$ is introduced. By Corollary~\CGref{cor:RH} the condition on $G$ and $A$ is then equivalent to the Riemann hypothesis. The introduction of $\zeta$ therefore does not reduce the condition to a weaker statement. The step that remains after it is the Riemann hypothesis itself. What $\zeta$ brings are the known theorems on its zeros, which become theorems on the solutions of $G=A$. The theorem of Hadamard and de la Vall\'ee Poussin excludes solutions on the line $\operatorname{Re}s=1$, which the methods of Sections~\CGref{sec:GA} and~\CGref{sec:B} do not do. The Riemann--von Mangoldt formula \CGeqref{eq:NT} counts the solutions with $0<\operatorname{Im}s\le T$, and the points $\tau_k$ of \CGeqref{eq:tauk} are defined from its main term. But \CGeqref{eq:NT} and the points $\tau_k$ concern only the imaginary parts of the solutions. They give no information on the real parts, and the points $\frac12+i\tau_k$ are not themselves solutions. The two routes therefore stop at the same place, the real parts of the solutions of $G=A$ in the strip $0<\operatorname{Re}s<1$. Without $\zeta$, the obstruction is that the structure proved so far is shared by kernels whose solutions lie off the line. With $\zeta$, the obstruction is the Riemann hypothesis.

It remains to say whether $G$ and $A$ intersect, and where. They do intersect, and this is not impossible to prove. On the real axis, $G(-1)=0$ by \CGeqref{eq:special} and $A(-1)=0$ by Proposition~\CGref{prop:A}, so $G(-1)=A(-1)$ without any use of $\zeta$, while $G(s)\ne A(s)$ for $0<s<2$ by Proposition~\CGref{prop:GAbound}. By \CGeqref{eq:GAzeta}, $G(s)=A(s)$ also at every odd integer $s\ge3$ and $s\le-3$, because $\cot(\frac{\pi s}2)=0$ there and $\zeta$ is finite at every $s\ne1$ \cite[Theorem~2.1, p.~13]{Titchmarsh1986}. Off the real axis, Corollary~\CGref{cor:RH} and Hardy's theorem that $\zeta$ has infinitely many zeros on the line $\operatorname{Re}s=\frac12$ \cite[\S10.2, p.~256]{Titchmarsh1986} show that $G=A$ at infinitely many points $\frac12+it$. So it is proved that $G$ and $A$ intersect infinitely often off the real axis, and on the line $\operatorname{Re}s=\frac12$ itself. Hardy's proof uses that $\zeta(\frac12+it)$ times an explicit factor is real, which is a consequence of the functional equation. The representations of Sections~\CGref{sec:GA} and~\CGref{sec:B} do not show a symmetry of $G-A$ under $s\mapsto1-s$, and this note does not prove the existence of a non-real solution of $G=A$ without $\zeta$. What is not proved is that $G$ and $A$ intersect off the real axis only on the line $\operatorname{Re}s=\frac12$. By Corollary~\CGref{cor:RH} this is the Riemann hypothesis. It is not known to be impossible to prove; it is open.
